
\begin{obereflection}
\textbf{Latihan 12.1}
1. Apa yang terjadi jika pengirim menggunakan nilai $k$ yang sama untuk dua pesan yang berbeda? Analisis risikonya terhadap keamanan pesan.
2. Jelaskan mengapa pemilihan $g$ sebagai akar primitif dari $p$ sangat penting dalam algoritma ElGamal.
3. Mengapa ukuran cipherteks ElGamal selalu dua kali lipat dari ukuran plainteksnya?
4. Diberikan parameter $p = 2357$, $g = 2$, dan kunci privat $x = 900$. Hitunglah
   kunci publik $y$, lalu enkripsikan $m = 1234$ dengan $k = 1000$. Periksa bahwa
   dekripsinya memang mengembalikan $1234$.
5. Cipherteks $(a_{1}, b_{1}) = (1439, 74)$ adalah hasil enkripsi $m_{1} = 700$
   pada modulus $p = 2273$. Bila penyerang memperoleh cipherteks lain
   $(a_{2}, b_{2}) = (1220, 1682)$ untuk pesan yang tidak diketahui, dapatkah ia
   memperoleh pesan kedua itu? Jelaskan langkah yang harus dilakukannya, atau
   jelaskan mengapa ia tidak dapat melakukannya.
6. Mengapa invers $\left(a^{x}\right)^{-1}$ dapat dihitung sebagai
   $a^{\,p-1-x} \bmod p$? Sebutkan teorema yang mendasarinya, dan syarat apa yang
   harus dipenuhi oleh $a$ agar teorema itu berlaku.
7. Jelaskan mengapa kunci privat $x = 1$ dan $x = p-1$ dikecualikan dari selang
   pemilihan kunci. Tunjukkan apa yang terjadi pada kunci publik $y$ dalam kedua
   kasus itu.
8. Misalkan $p = 2273$ dan pesan yang dikirim adalah rangkaian delapan digit
   $07001114$. Bandingkan pemecahan menjadi blok empat digit dengan pemecahan
   menjadi blok tiga digit. Berapa blok yang dihasilkan masing-masing cara, dan
   masalah apa yang muncul pada pemecahan tiga digit?
\end{obereflection}

\begin{obereflection}
\textbf{Latihan 12.2}
1. Sistem ElGamal memakai $k$ baru untuk setiap blok pesan, sehingga satu blok
   memerlukan dua perpangkatkan modular. Bandingkan biaya itu dengan RSA, yang
   hanya memerlukan satu perpangkatkan per blok, dan jelaskan mengapa ElGamal
   tetap dipandang layak dipakai.
2. Sebuah berkas multimedia berukuran 1 MB hendak dikirim dengan aman. Mengapa
   tidak tepat mengenkripsinya langsung dengan ElGamal? Rancanglah skema
   kriptografi hibrida yang memakai ElGamal untuk membungkus kunci sesi dan AES
   untuk isinya, lalu sebutkan besaran apa saja yang melewati saluran.
3. Tunjukkan bahwa ElGamal bersifat homomorfik multiplikatif. Perlihatkan
   langkahnya untuk $m_{1} = 700$ dan $m_{2} = 1114$ pada modulus $2273$, lalu
   jelaskan mengapa sifat itu berbahaya bagi keutuhan pesan.
4. OpenPGP pernah mencantumkan ElGamal sebagai algoritma kunci publik nomor 16
   dengan keterangan \textit{encrypt-only}. Jelaskan mengapa keterangan itu
   diberikan, dan mengapa RFC~9580 (2024) kemudian melarang pembangkitan kunci
   ElGamal yang baru.
5. Bandingkan ElGamal dengan protokol Diffie--Hellman dari sudut kebutuhan
   interaksi. Dalam keadaan apa keunggulan ElGamal itu menjadi menentukan, dan
   dalam keadaan apa keunggulan itu tidak berguna?
6. Jelaskan mengapa algoritma Shor mengancam baik RSA maupun ElGamal, sekalipun
   keduanya bersandar pada persoalan yang berbeda. Apa yang harus dilakukan
   sebuah lembaga yang hari ini menyimpan data rahasia jangka panjang?
\end{obereflection}
